\documentclass[11pt,a4paper]{ctexart}

\usepackage[margin=2.2cm,top=1.8cm,bottom=1.8cm]{geometry}
\usepackage{booktabs}
\usepackage{array}
\usepackage{listings}
\usepackage{xcolor}

\lstset{
  basicstyle=\ttfamily\normalsize,
  columns=flexible,
  breaklines=true,
  frame=single,
  framesep=3pt,
  xleftmargin=4pt,
  aboveskip=5pt,
  belowskip=5pt,
  keepspaces=true,
}

\setlength{\parindent}{2em}
\setlength{\parskip}{3pt}
\linespread{1.12}
\renewcommand{\arraystretch}{1.15}
\emergencystretch=1.5em
\newcommand{\code}[1]{\texttt{#1}}

\begin{document}

\begin{center}
{\large\bfseries 实验二：C{-}{-} 语义分析与类型检查}\\[3pt]
张苏楠\quad 241880199
\end{center}

\textbf{工具版本}：gcc \textbf{11.4.0}、GNU flex \textbf{2.6.4}、GNU bison
\textbf{3.8.2}、GNU make \textbf{4.3}（Ubuntu 22.04）

\section*{一、实现的功能，以及如何实现}

{\normalsize
\setlength{\parindent}{0pt}
\begin{tabular}{@{}>{\raggedright\arraybackslash}m{2.9cm}m{13.1cm}@{}}
\hline
\textbf{功能} & \textbf{如何实现} \\
\hline
19 个错误码\par（必做 17 + 选做 2.1 的 18、19） & 前端仍是实验一的 Flex + Bison（文法只加了一条产生式），\textbf{语义分析完全手写}：在实验一建好的语法树上做一遍后序遍历。三个模块单向依赖 —— \code{semantic.c}（遍历与 19 种检查）$\to$ \code{symtab.c}（符号表）$\to$ \code{type.c}（类型表示与等价）。\\
\hline
错误 1/2 的\textbf{连锁抑制} & 名字查不到时返回 \code{UNKNOWN} 单例并报 1/2；\code{typeEqual(UNKNOWN, 任何)} 为真，各检查点遇 \code{UNKNOWN} 一律静默。否则 \code{j = i + 1;}（\code{j} 未定义）会附带一条类型错 —— 要求书只考察「较本质」的那条。\\
\hline
嵌套作用域\par（选做 2.2） & 符号表是散列 + 每层一条链表；查表沿链由内向外，返回\textbf{第一个 \code{depth} $\le$ 当前深度}的项 —— 这就是「遮蔽」。\code{scopeEnter}/\code{scopeExit} 全项目只有两处调用点（函数体、语句块），形参与函数体同层。\\
\hline
函数声明\par（选做 2.1） & 声明的形参收集与定义的走\textbf{同一条}路径，比较（返回类型, 形参数, 形参类型表）；文件末尾遍历全局链表报「声明了但从未定义」（错误 18），行号取\textbf{首次}声明处。\\
\hline
结构等价\par（选做 2.3） & 忽略域名、按位置比域的类型，域数量必须相同；数组整体作为一个（基类型, 维数）单位比较，\textbf{不展开}（为什么必须如此，见第三节）。\\
\hline
输出格式 & 词法/语法错仍报 \code{A}/\code{B}；\textbf{完全无错时不输出任何内容}（不再打印语法树）。语义报告走 \code{reportSemanticError}，它维护一个\textbf{已打印过的行号集合}，重复行直接丢弃（判据无条件，不依赖「相邻」）；实验一的 \code{A}/\code{B} 路径原样保留，两条路径的去重状态互不干扰。\\
\hline
\end{tabular}}

\section*{二、如何编译}

\begin{lstlisting}[language=bash]
sudo apt-get install -y flex bison gcc make
make clean && make          # 生成 ./parser（同时生成同内容的 ./cc）
./parser <文件>              # 无错时不输出任何内容；有错时打印错误信息
\end{lstlisting}

编译与实验一完全相同：仓库根目录执行 \code{make} 即可，除 \code{src/} 外不依赖任何
目录。\textbf{请勿使用 \code{make -j}} —— \code{syntax.tab.c syntax.tab.h} 是一条
「一个 recipe、两个目标」的多目标规则，并发下会被当成两条独立规则、同时跑两次 bison
并写同一个文件。

关于 Bison 版本：文法不使用 3.1 以后才引入的写法（不用
\code{\%define parse.error detailed}、不用 \code{\%union}），因此 \textbf{3.0.4 及
以上各版本均可编译}；实测 3.0.4 / 3.5.1 / 3.8.2 三个版本下，冲突 state 的编号与状态项
\textbf{逐字相同} —— 编号只取决于文法，与版本无关。

\section*{三、亮点 1：把「不展开数组」从要求书的例子里反推出来}

要求书 2.3 给了一对\textbf{结构等价}的例子（域名不同：\code{\{int x; float y;\}}
$\equiv$ \code{\{int y; float z;\}}），又给了一对\textbf{不等价}的例子。后者初看
令人费解：那两个结构体「看起来挺像」，凭什么不等价？把最自然的那种实现 —— 结构等价
$=$ 把各域的类型\textbf{展开}成一个序列再逐位比较 —— 代入就明白了：

\begin{lstlisting}[language=C]
A = struct A { int a; struct { float f; int i; } b[10]; }  -> int, [float,int]x10  = 21 项
B = struct B { struct { int i; float f; } b[10]; int b; }  -> [int,float]x10, int  = 21 项
\end{lstlisting}

展开后 A 与 B \textbf{逐位完全相同}（\code{int==int}、\code{float==float}、\dots），
于是会把要求书明确判定为\textbf{不等价}的一对判成等价。所以那句话不是修辞，它是在说
「别展开」：A 的第 0 个域是 \code{int}、B 的第 0 个域是 \code{array}，\textbf{种类就
已经不同}，立刻判出不等价。代码上它只是让结构等价逐个比较域的类型，而数组域整体走
\code{typeEqual} 的数组分支（把整个数组当作一个不可分割的（基类型, 维数）单位）。

值得注意的是，这对「不等价」的例子同时是「忽略域名」的独立佐证（它的两个结构体域名正是
不同的），而「忽略域名」的直接佐证又只能来自那对「等价」的例子。\textbf{推理到这里才
闭合}：两条规则互为对方的证据，任何一条单独看都像是凭空规定。

\section*{四、亮点 2：一个在语义层\textbf{观测不到}的 bug，决定了模块边界}

\code{int a[10][2]} 的 \code{VarDec} 树是左递归的：
\code{VarDec(VarDec(a) LB 10 RB) LB 2 RB} —— \textbf{最外层结点带的是最后一个方括号}。
所以类型必须\textbf{从右往左}折叠；若照收集顺序从左折叠，得到的是
\code{array(2, array(10, int))}：\textbf{维度顺序反了}。

这个 bug 在本实验里\textbf{没有任何症状}：数组等价\textbf{忽略长度}，反了的数组与
任何同维数组比较仍然「相等」；\textbf{维数没变}，于是错误 10（对非数组用 \code{[]}）
的触发条件不变；错误 12 只看\textbf{下标}的类型，与数组本身是什么类型无关。也就是说，
17 种检查\textbf{一条都发现不了它} —— 整个语义分析阶段无从暴露它。

唯一能钉住它的办法，是让类型模块\textbf{脱离词法、语法、符号表独立可测}，并且断言
\textbf{具体的 \code{size} 值}（只查维数同样抓不到）。这就是这条模块边界的由来：
\code{type.c} \textbf{不依赖任何东西}（连符号表都不依赖），不是为了分层好看，而是因为
\textbf{只有它能看见这个 bug}。一句话：\textbf{判别力的边界决定模块的边界} —— 一个错误
若在高层没有任何可观测量，就只能靠把低层做成可独立测试的纯模块来发现。

\section*{五、亮点 3：把「要求书没规定」的十条做成可单点翻转的开关}

要求书与指导书对十种情形\textbf{都没有写}：变量名与函数名冲突报哪个码？\code{if} 的
条件不是 \code{int} 报错误 7 还是别的？形参与函数体算不算同一个作用域？\code{a[b]}
（\code{b} 是数组）报 12 还是 7？同一行出现两条错误留哪一条？这些\textbf{推不出来}，
两份文档里都没有依据，只能自己选定一种行为。

既然依据无从查出，目标就不该是「猜得准」，而是\textbf{改得起}。为此每条选定行为都必须
满足三点，缺一不可。

\begin{enumerate}
\item \textbf{行为收敛在代码里的一个位置}，并注明「这是第 N 条，要求书未规定」
\allowbreak：\code{symConflictCode()}（1）\allowbreak、\code{checkCondition()}（2）\allowbreak、
\code{analyzeExtDef} 的 \code{ED\_FUNC} 分支里\textbf{那一次} \code{scopeEnter()} 调用
（3，必须落在登记函数名之后、登记形参之前）\allowbreak、
\code{analyzeStructSpecifier} 里定义结构体那条路径上的 \code{symEnter()}（4，用的是
\textbf{当前}作用域深度）\allowbreak、\code{analyzeStructSpecifier} 里
\code{nchild == 2} 的那个分支（5）\allowbreak、\code{analyzeSubscript}（6）\allowbreak、
\code{symEnter} 的返回值语义（7）\allowbreak、
\code{checkable()} 毒化开关（8）\allowbreak、\code{report.c} 的同行去重（9）。
做不到这一点，它就会退化成散在代码里的偶然行为，改一条要通读全文。
\item \textbf{各配一个专属用例}钉住当前行为（\code{Test/lab2/dec01}--\code{dec17}）——
否则改动其中一条时，没有任何东西会变红。
\item \textbf{一张从症状反查决策的对照表}：
\end{enumerate}

{\normalsize
\setlength{\parindent}{0pt}
\begin{tabular}{@{}m{8.2cm}@{\hspace{8pt}}m{7.4cm}@{}}
\hline
\textbf{观测到的症状} & \textbf{最可能的决策} \\
\hline
未定义变量的表达式被多报了一条类型错 & 8（\code{UNKNOWN} 毒化）\\
\hline
同一行出现两条错误 & 9（同行只留最先检出的）\\
\hline
某个跨函数的同名局部变量被报重复定义 & 3 或 4（作用域划在哪一层）\\
\hline
\code{if (3.7)} 被报错 / 该报错却没报 & 2（条件必须是 \code{int}）\\
\hline
数组下标里放了个数组，报的码是 7 而不是 12 & 6（专项优先于通则）\\
\hline
结构体与变量同名时报的码不对 & 1（单一名字空间，按后定义者报码）\\
\hline
重复定义的变量，后续引用被解析到了后一个 & 7（保留最早那个）\\
\hline
\code{struct Undefined;}（无变量）没报 17 & 5（无变量也报）\\
\hline
使用了文件后面才定义的变量/函数 & 10（定义必须先于使用）\\
\hline
\end{tabular}}

第 10 条要单独说：它的「开关」不是一行代码，而是\textbf{遍历结构本身}（单遍、不做
预扫描，于是定义必须先于使用），也是十条里唯一「改起来代价不对称」的一条 —— 若要改成
「全文件任意位置定义即可」，得先加一趟预扫描收集全局名，而不是改一个 \code{if}。正因
如此，它最需要这张表和那条用例。这张表才是整套做法的落点：十种情形里任何一种若被判定
需要改，它是\textbf{从症状直接定位到那一个开关}的东西 —— 把未知收敛成已知，靠的不是
猜得准，而是\textbf{选定过的行为可查、可改、可验证}。

同一套思路也用在期望文件上：每个官方样例配两个期望文件 —— \code{.exp} 钉我们自己的
完整文本、\code{.pairs} 钉官方的「错误类型 + 行号」序列（要求书明说这是唯一的判据），
两者分开，才能把「改了说明文字」与「错误类型/行号错了」分开暴露。

最后一条原则贯穿全篇：\textbf{每一条检查都要自证有判别力} —— 把它改坏，必须有用例变红。
照这条原则反向验证时，确实发现并修掉了两处「判据看似存在、实则没有判别力」：一处是
\code{scopeExit()} 的摘链动作\textbf{当时没有任何检查看着}（原有检查点全在其后，此时
\code{curDepth == 0}，\code{symLookup} 的深度守卫自己就把残留过滤掉了，于是「退层不摘
散列链」这个改动\textbf{无人察觉}）；另一处是 \code{UNKNOWN} 毒化的探针\textbf{两个候选
错误全落在同一行}，被同行去重抑制，把毒化整个关掉也不变红 —— 拆到不同行才恢复判别力。
没有判别力的判据会和缺陷一起交付，所以这条不是自我要求，是必要条件。

\end{document}
